\section*{章末問題}
\begin{enumerate}[label=\textbf{\thechapter.\arabic*},leftmargin=*]
 \item\label{問:区分求積分子が添字}
 次の極限を求めよ。
 \[
  \lim_{n\to\infty}\sum_{k=1}^{n}\frac{k}{n^2+k^2}.
 \]
 \item\label{問:区分求積根号}
 次の極限を求めよ。
 \[
  \lim_{n\to\infty}\cfrac1{n^{\cfrac32}}
      \sum_{k=1}^{n}\sqrt{n+k}.
 \]
 \item\label{問:逆平方和の残り}
 自然数$N>n\geqq1$に対して
 \[
  \frac1{n+1}-\frac1{N+1}
  \leqq\sum_{k=n+1}^{N}\frac1{k^2}
  \leqq\frac1n-\frac1N
 \]
 を示せ。これを使って$\sum_{k=1}^{\infty}\frac1{k^2}$が収束し，その第$n$項より後の和$R_n$が
 $\frac1{n+1}\leqq R_n\leqq\frac1n$を満たすことを示せ。
 \item\label{問:積分制約と二乗}
 $f$は$[0,1]$で連続で，$\int_0^1xf(x)\dd x=1$を満たすとする。
 $\int_0^1f(x)^2\dd x$の最小値を求め，それを与える関数を求めよ。
 \item\label{問:対数の積分漸化式}
 非負整数$n$に対して$K_n=\int_1^{\e}(\log x)^n\dd x$とおく。
 $K_0$と，$n\geqq1$における漸化式を求め，$K_3$を計算せよ。
 \item\label{問:ウォリス積の誤差}
 $P_m=\prod_{k=1}^{m}\frac{4k^2}{4k^2-1}$とする。
 \[
 0\leqq\frac{\pi}{2}-P_m\leqq\frac{\pi}{4m+4}
 \]
 を示せ。また$\lim_{m\to\infty}\sqrt m\int_0^{\frac{\pi}{2}}\cos^{2m}x\dd x$を求めよ。
 \item\label{問:整数ベータ二項係数}
 非負整数$m,n$に対して
 \[
 \int_0^1\binom{m+n}{m}x^m(1-x)^n\dd x=\frac1{m+n+1}
 \]
 を示せ。次に$\int_0^1x^2(1-x)^4\dd x$を求めよ。
 \item\label{問:半整数ベータ}
 広義積分$\int_0^1\sqrt{\frac{x}{1-x}}\dd x$の収束を両端で確かめ，ベータ関数の漸化式を使って値を求めよ。さらに$x=\sin^2t$による置換でも検算せよ。
\end{enumerate}

\smallskip\noindent 解答は\pageref{解答:第9章}ページ。
